\section{Implementation}
\label{sec:implementation}

We briefly describe how \system realizes the oblivious primitives of
\cref{sub:obl-primitives}. \cmpset is implemented with branchless
conditional-move instructions; rows are laid out in 64-byte blocks, so
obliviously selecting or swapping an entire row compiles to a single SIMD
masked-blend instruction. All higher-level operators are built on this
primitive: \textsc{OSort} is a bitonic sorting
network~\cite{batcher1968sorting}, \textsc{OCompact} is
ORCompact~\cite{sasy2022fast}, and \oht follows
H2O2RAM~\cite{zheng2025h2o2ram}, employing a planner that selects among
linear scan, bucket hash, stashless Cuckoo hash, and two-tier hash based on
the public table size and lookup count. Notably, \textsc{Dedup} and
\textsc{Redup} require no generic sort-and-compaction: since edge lists are
stored sorted by endpoint, each is a single linear pass that uses \cmpset
to obliviously suppress a duplicate row (replacing it with a dummy) or
restore a suppressed one (copying the preceding row forward). Finally,
\system parallelizes all of its operators, but every parallel schedule is
derived exclusively from public parameters --- the fixed structure of the
sorting and compaction networks, the chunking of linear passes by public
table sizes, and the concurrent execution of the source and destination
branches of \onehop --- so parallel execution reveals nothing beyond the
sequential trace.
